%% Cover letter for the Cognitive Neurodynamics submission.
%%   pdflatex cover-letter.tex && cp cover-letter.pdf ../paper-submission/
\documentclass[11pt,a4paper]{article}
\usepackage[T1]{fontenc}
\usepackage[utf8]{inputenc}
\usepackage{newtxtext}
\usepackage[a4paper,margin=2.5cm]{geometry}
\usepackage[hidelinks]{hyperref}

\pagestyle{empty}
\setlength{\parindent}{0pt}
\setlength{\parskip}{1em}

\begin{document}

3 August 2026

Dear Editor,

We wish to submit ``Adult neurogenesis promotes pattern separation and sparsity beyond
the dentate gyrus in a spiking DG--CA3 model'' for consideration as a Research Article.

Existing models of adult neurogenesis and pattern separation stop at the dentate gyrus.
Using a spiking dentate gyrus--CA3 network built from Hippocampome.org and extended with
immature granule cells, we show that neurogenesis makes activity sparser and improves
pattern separation in CA3 as well, reproducing \textit{in vivo} observations (McHugh et
al.\ 2022) from the known circuitry alone. This extends work the journal has published
(Kim and Lim 2024a, b; Yang et al.\ 2025).

The manuscript is original, is not under consideration elsewhere, and all authors have
approved it. It derives from the first author's master's thesis at the Federal University
of Para\'iba. Code and data: \url{https://github.com/MarlonVCendron/neurogenesis}. No
funding was received, the authors declare no competing interests, and the study is
entirely computational.

Sincerely,

\vspace{1.5em}

Marlon Valm\'orbida Cendron (corresponding author)\\
Federal University of Para\'iba, Jo\~ao Pessoa, Brazil\\
\href{mailto:marlonvcendron@gmail.com}{marlonvcendron@gmail.com}

\end{document}
